\section{SM pipelines}
\label{sec:sm}

\subsection{The core table}

The four-cycle class \citep[tab.~4.1]{jia2019turing} reproduces:
\texttt{FFMA}, \texttt{FADD}, \texttt{FMUL}, \texttt{LOP3}, \texttt{SHF},
\texttt{SEL}, and the derived \texttt{IADD3}, \texttt{ISETP}, and
\texttt{FSETP} all measure between 3.96 and 4.13 cycles. \texttt{POPC}
and \texttt{FLO} measure exactly 15, \texttt{HFMA2} 6.19, both on prior.
Three deviations publish: \texttt{IMAD} at 4.11 cycles against a prior of
5 ($-17.8\%$), \texttt{DADD} at 44.0 against ${\sim}48$, and
\texttt{DFMA} at 48.0 against ${\sim}54$. Rows with no published prior
are contributions: \texttt{PRMT} at 4.13 cycles, \texttt{IDP.4A} at 4.11
cycles and the full 2.0 warp-instructions per SM per cycle, and the
emulated \texttt{u32} divide at 55.5 cycles per sequence. Throughput
saturates at 2.0 warp-instructions per SM per cycle for both the FP32 and
INT32 classes (64 lanes of each per SM), with the FFMA anchor
landing within 0.005\% of the geometric 2.0. \Cref{tab:core} collects
these rows. The repository table carries the remaining variants and the
provenance columns.

\begin{table}[t]
\centering
\caption{Core instruction rows. Latency by dependent \texttt{clock64}
chain, throughput at the peak of the warps-per-SM sweep, pipe bindings by
contention probe (\cref{sec:methodology}). The pipe \emph{own} denotes a
dedicated unit. Starred rows derive from pair constructions. Latency priors from
Jia et al.\ \citep{jia2019turing}.}
\label{tab:core}
\small
% generated by tools/mk_paper_tables.py -- do not edit
\begin{tabular}{@{}l l r r r r@{}}
\toprule
 & & \multicolumn{3}{c}{latency (cycles)} & throughput \\
\cmidrule(lr){3-5}
op & pipe & meas. & prior & dev.\,\% & (wi/SM/clk) \\
\midrule
\texttt{FFMA} & fma & 4.1 & 4 & $+1.8$ & 2.00 \\
\texttt{FADD} & fma & 4.1 & 4 & $+2.9$ & 2.00 \\
\texttt{FMUL} & fma & 4.1 & 4 & $+2.9$ & 2.00 \\
\addlinespace
\texttt{IMAD} & fma & 4.1 & 5 & $-17.8$ & 2.00 \\
\texttt{IADD3*} & -- & 4.0 & 4 & $-1.0$ & 1.98 \\
\texttt{LOP3} & alu & 4.1 & 4 & $+2.7$ & 1.97 \\
\texttt{SHF} & alu & 4.1 & 4 & $+2.7$ & 1.97 \\
\texttt{SEL} & alu & 4.1 & 4 & $+2.5$ & 1.95 \\
\texttt{PRMT} & alu & 4.1 & -- & -- & 1.97 \\
\texttt{ISETP*} & -- & 4.0 & 4 & $-0.8$ & 1.98 \\
\texttt{IDP.4A.S8} & fma & 4.1 & -- & -- & 2.00 \\
\addlinespace
\texttt{HFMA2} & own & 6.2 & 6 & $+3.1$ & 1.96 \\
\texttt{DADD} & own & 44.0 & 48 & $-8.3$ & 0.05 \\
\texttt{DFMA} & own & 48.0 & 54 & $-11.1$ & 0.05 \\
\addlinespace
\texttt{POPC} & own & 15.0 & 15 & $+0.0$ & 0.50 \\
\texttt{FLO} & own & 15.0 & 15 & $+0.0$ & 0.50 \\
\addlinespace
\texttt{MUFU.EX2} & -- & 15.0 & 15 & $+0.0$ & 0.50 \\
\texttt{MUFU.RSQ} & -- & 15.0 & -- & -- & 0.50 \\
\texttt{MUFU.SIN} & -- & 21.0 & -- & -- & 0.50 \\
\addlinespace
\texttt{SHFL.IDX} & -- & 25.0 & -- & -- & 0.50 \\
\bottomrule
\end{tabular}

\end{table}

Some rows are measurable only as pairs from PTX: a predicate cannot be
chained, so \texttt{ISETP} latency derives from a compare-plus-select
chain minus the select chain. A pure dependent add chain cannot be
pinned at all (every guard was strength-reduced), so \texttt{IADD3}
derives from an add-xor pair minus the \texttt{LOP3} row. Each derived
row names its construction.

\subsection{Pipe bindings}

Contention probes (\cref{sec:methodology}) bind each operation to its
issue pipe by mixing it 50/50 with a known reference and comparing the
combined rate against same-pipe (harmonic) and distinct-pipe (maximum,
capped at the 4-per-SM-cycle issue rate) predictions. The probe
self-tests pass: FFMA against itself reads same-pipe, FFMA against
LOP3 reads distinct. The map follows. \texttt{IMAD} issues on the
fma pipe, confirming the Turing arrangement by measurement.
\texttt{SHF}, \texttt{SEL}, and \texttt{PRMT} sit on the alu pipe.
\texttt{POPC} and \texttt{FLO} run on a quarter-rate unit that further
\emph{partially} contends with the load path (a mixed POPC-and-load
stream runs at 0.68 of the distinct-pipe prediction). And,
consequentially for quantised inference kernels, \textbf{\texttt{IDP.4A}
issues on the fma pipe}, competing with \texttt{FFMA} for issue slots
rather than riding the integer pipe.

\subsection{Tensor cores}

The m16n8k8 \texttt{HMMA} with f16 accumulate measures 0.50000
warp-instructions per SM per cycle, exactly the whitepaper-derived 512
FMA per SM per cycle \citep{nvidia2018turing}. \texttt{IMMA} m8n8k16
likewise lands exactly on its derived 1.0. The f32-accumulate prior of
half rate is \emph{refuted on this part}: measured 0.49996, identical to
f16 accumulate. The half-rate figure describes GeForce-positioned TU102s.
The Quadro runs full-rate f32 accumulation: a market-segmentation
fact rather than a microarchitectural one, and worth a row precisely
because the whitepaper does not say it. A dependent accumulate chain puts
\texttt{HMMA} latency at 14.0 cycles, and a contention probe confirms
that \texttt{HFMA2} shares the tensor unit (mixed rate 0.76 against a
shared-unit prediction of 0.80): FP16 math runs through the tensor
cores, measured rather than inferred. The operand-staging load,
\texttt{LDSM} (\texttt{ldmatrix}), chains at 30 cycles through an
address perturbation and peaks at 0.125 warp-instructions per SM per
cycle in its four-tile form. That is 512 bytes per instruction, exactly
the 64-byte-per-cycle unified-datapath ceiling of \cref{sec:memory}:
feeding the tensor cores is bandwidth-bound on the same path as every
other shared-memory access. \Cref{tab:tensor} summarises.

\begin{table}[t]
\centering
\caption{Tensor-core rows. The f32-accumulate rate equals the
f16-accumulate rate on this Quadro-positioned part. The half-rate figure
in circulation describes GeForce-positioned TU102s.}
\label{tab:tensor}
\small
% generated by tools/mk_paper_tables.py -- do not edit
\begin{tabular}{@{}l r r r@{}}
\toprule
op & tput (wi/SM/clk) & prior & dev.\,\% \\
\midrule
\texttt{HMMA.1688}, f16 accumulate & 0.500 & -- & -- \\
\texttt{HMMA.1688}, f32 accumulate & 0.500 & -- & -- \\
\texttt{IMMA.8816}, s8 & 1.000 & 1.0 & $-0.0$ \\
\addlinespace
\texttt{HMMA.1688} dependent latency & \multicolumn{3}{r}{14.0 cycles} \\
\bottomrule
\end{tabular}

\end{table}

\subsection{SFU, shuffle, divergence}

The remaining \texttt{MUFU} set sits at 15 cycles (\texttt{RSQ},
\texttt{LG2}) with \texttt{SIN} and \texttt{COS} as 21-cycle pairs
carrying an \texttt{FMUL} range scale, all at the quarter rate.
\texttt{RCP} derives at ${\sim}17$ cycles after the approximate-algebra
deletions of \cref{sec:sass}. \texttt{SHFL.IDX} chains at 25 cycles and
0.5 per cycle. A ballot (\texttt{VOTE.ANY}) chains at 16.2 cycles and
sustains 0.96 per cycle across independent chains, with the predicate
lane-shifted to keep the chain off the uniform datapath. A barrier in a
192-thread block costs 32.3 cycles whether or not a second 192-thread
block is resident on the same SM: the barrier unit serves two
concurrent CTAs without serialising them. Branch divergence is linear:
a $k$-way divergent switch costs close to $k$ times the uniform path
(16.7 cycles to 528 for 1-way to 32-way, within 5\% of $k\times$
throughout), and the predicated equivalent
costs precisely the uniform 16.7. If-conversion is free at this
shape, with reconvergence visible as \texttt{BSSY}/\texttt{BSYNC} in the
SASS.
